\begin{abstract}
LLM-powered analytics answers group-level decision queries, such as HAVING thresholds, rankings, and top-$k$ selections over LLM-judged attributes, with one model call per tuple.
Sampling reduces this cost, but a sampled decision needs a guarantee that survives continuous monitoring.
Confidence sequences can stop once a threshold claim is settled, but they bound each group's mean on its own and ignore the proxy scores the system already holds; proxy-based estimators use those scores but certify estimates at a fixed sample size, not decisions under continuous monitoring.
We present \sys{}, a decision layer that compiles HAVING, ORDER BY, and Top-$k$ queries into a handful of boundary hypotheses and tests each with a betting e-process, which stays valid at any stopping time by Ville's inequality.
It bets on proxy-rectified values, which keep every group mean exact for any proxy while removing the variance that the proxy explains.
Finite-population and paired forms of the tests handle sampling without replacement and groups of very different sizes.
We prove that \sys{} returns an incorrect answer with probability at most $\alpha$ under any stopping rule and any proxy.
On \numNQ{} queries over three datasets, including 122 thousand reviews and 309 thousand paper abstracts judged by a real LLM, \sys{} returned no incorrect answer in \numNEXEC{} executions.
It used \numRWSR$\times$ fewer LLM calls than Evergreen-style confidence sequences (up to \numRWSRMAX$\times$ on large groups) and \numRHS--\numRBSP$\times$ fewer than IFOCUS- and BlazeIt-style bounds, and about as many as uncertified normal-approximation stopping, which returned wrong answers in up to 13\% of a query's executions.
With two LLMs in the loop, latency followed the number of calls, and HAVING queries ran \numLiveHavingMin$\times$ and \numLiveHavingMax$\times$ faster than full evaluation.
\end{abstract}
